\documentclass[11pt,a4paper]{article}
\usepackage{i2c_report_style}

\begin{document}

\makekyureportheader{LABORATORY REPORT: TASK 3}{Character LCD1602 Interfacing via PCF8574 I2C Backpack}

\begin{objectivebox}
The objective of this laboratory task is to interface an alphanumeric HD44780-compatible $16\times 2$ character LCD over the $\text{I}^2\text{C}$ bus using a PCF8574 8-bit expander backpack. Students analyze 4-bit nibble packetization, evaluate address mismatch fault handling, implement dynamic horizontal text scrolling, and calculate wire-level transmission throughput and display refresh bandwidth.
\end{objectivebox}

% =========================================================================
\section{Hardware Architecture \& Backpack Mapping}
% =========================================================================
The standard Hitachi HD44780 LCD controller requires a minimum of 6 control/data lines in 4-bit parallel mode ($\text{RS}, \text{E}, \text{D}_4, \text{D}_5, \text{D}_6, \text{D}_7$). The "1602 IIC" backpack employs a PCF8574 remote 8-bit I/O expander to bridge the two-wire $\text{I}^2\text{C}$ bus to the parallel LCD interface pins.

\begin{table}[h!]
\centering
\small
\caption{PCF8574 Output Pin to HD44780 LCD Signal Mapping}
\vspace{4pt}
\begin{tabularx}{\textwidth}{l c c X}
\toprule
\textbf{PCF8574 Port} & \textbf{Bit Mask} & \textbf{LCD Signal} & \textbf{Functional Description} \\
\midrule
$P_0$ & \texttt{0x01} & \textbf{RS} (Register Select) & \texttt{0} = Command Register, \texttt{1} = Data (Character) Register \\
$P_1$ & \texttt{0x02} & \textbf{RW} (Read/Write) & Held permanently at \texttt{0} (Write mode only) \\
$P_2$ & \texttt{0x04} & \textbf{E} (Enable) & High-to-Low transition ($E \downarrow$) latches data on $D_4\text{--}D_7$ \\
$P_3$ & \texttt{0x08} & \textbf{Backlight} & Controls NPN transistor for LED backlight (\texttt{1} = ON, \texttt{0} = OFF) \\
$P_4$ & \texttt{0x10} & \textbf{D4} & LCD Data Bit 4 \\
$P_5$ & \texttt{0x20} & \textbf{D5} & LCD Data Bit 5 \\
$P_6$ & \texttt{0x40} & \textbf{D6} & LCD Data Bit 6 \\
$P_7$ & \texttt{0x80} & \textbf{D7} & LCD Data Bit 7 (also Busy Flag $BF$ in read mode) \\
\bottomrule
\end{tabularx}
\end{table}

\begin{theorybox}
\textbf{Enable Latch Timing Physics:} To transfer one byte to the LCD in 4-bit mode, the byte is partitioned into two 4-bit nibbles. Each nibble is sent twice: first with $E=1$ (\texttt{0x04}), then with $E=0$. The falling edge of $E$ latches the nibble into the LCD. Because each $\text{I}^2\text{C}$ byte transfer over a $100\text{ kHz}$ bus takes $90\,\mu\text{s}$, it inherently satisfies the HD44780 minimum enable pulse width ($t_{w(E)} \ge 450\text{ ns}$) without requiring microsecond delay instructions.
\end{theorybox}

% =========================================================================
\section{Firmware Implementation}
% =========================================================================

\subsection{Platform 1: STM32 HAL Low-Level Driver (main.c)}
In the STM32 HAL implementation, we construct an atomic 4-byte buffer for every character: $[ \text{High Nibble }|\; E, \; \text{High Nibble}, \; \text{Low Nibble }|\; E, \; \text{Low Nibble} ]$ and transmit it in a single $\text{I}^2\text{C}$ transaction.

\begin{lstlisting}[language=C, caption={STM32 HAL PCF8574 LCD Driver with Scrolling Support}]
#define LCD_ADDR (0x27 << 1) // 7-bit 0x27 shifted left by 1
#define LCD_RS   0x01
#define LCD_EN   0x04
#define LCD_BL   0x08

static void lcd_nibble(uint8_t nib, uint8_t rs) {
    uint8_t v = (uint8_t)((nib << 4) | LCD_BL | (rs ? LCD_RS : 0));
    uint8_t d[2] = { (uint8_t)(v | LCD_EN), v };
    HAL_I2C_Master_Transmit(&hi2c1, LCD_ADDR, d, 2, 50);
}

static void lcd_byte(uint8_t b, uint8_t rs) {
    uint8_t hi = (uint8_t)((b & 0xF0) | LCD_BL | (rs ? LCD_RS : 0));
    uint8_t lo = (uint8_t)(((b << 4) & 0xF0) | LCD_BL | (rs ? LCD_RS : 0));
    uint8_t d[4] = { (uint8_t)(hi | LCD_EN), hi, (uint8_t)(lo | LCD_EN), lo };
    HAL_I2C_Master_Transmit(&hi2c1, LCD_ADDR, d, 4, 50);
}

void lcd_init(void) {
    HAL_Delay(50);
    lcd_nibble(0x03, 0); HAL_Delay(5);
    lcd_nibble(0x03, 0); HAL_Delay(1);
    lcd_nibble(0x03, 0); HAL_Delay(1);
    lcd_nibble(0x02, 0); HAL_Delay(1); // Force 4-bit mode
    lcd_byte(0x28, 0); // 4-bit mode, 2-line display, 5x8 font
    lcd_byte(0x0C, 0); // Display ON, Cursor OFF, Blink OFF
    lcd_byte(0x01, 0); HAL_Delay(2); // Clear display
    lcd_byte(0x06, 0); // Entry mode: increment cursor
}

void lcd_set_cursor(uint8_t row, uint8_t col) {
    lcd_byte((uint8_t)(0x80 | (row ? 0x40 : 0x00) | col), 0);
}

void lcd_print(const char *s) {
    while (*s) lcd_byte((uint8_t)*s++, 1);
}

/* Horizontal scrolling implementation for strings > 16 characters */
void lcd_scroll_message(uint8_t row, const char *msg, uint16_t delay_ms) {
    int len = strlen(msg);
    char window[17];
    for (int i = 0; i <= len - 16; i++) {
        strncpy(window, msg + i, 16);
        window[16] = '\0';
        lcd_set_cursor(row, 0);
        lcd_print(window);
        HAL_Delay(delay_ms);
    }
}
\end{lstlisting}

\subsection{Platform 2: Arduino Uno Sketch (Wire + LiquidCrystal\_I2C)}
\begin{lstlisting}[language=C++, caption={Arduino Uno LCD Program with Dynamic Scrolling}]
#include <Wire.h>
#include <LiquidCrystal_I2C.h>

LiquidCrystal_I2C lcd(0x27, 16, 2);

const char line1_msg[] = "Hello, Uganda";
const char scroll_msg[] = "    Kyambogo University Dept of EEE - TETE 3102 Microcontrollers    ";

void setup() {
  Serial.begin(115200);
  lcd.init();
  lcd.backlight();
  lcd.setCursor(0, 0);
  lcd.print(line1_msg);
}

void loop() {
  int len = strlen(scroll_msg);
  char window[17];
  for (int i = 0; i <= len - 16; i++) {
    strncpy(window, scroll_msg + i, 16);
    window[16] = '\0';
    lcd.setCursor(0, 1);
    lcd.print(window);
    delay(300);
  }
}
\end{lstlisting}

% =========================================================================
\section{Fault Injection: Address Mismatch Investigation}
% =========================================================================
In this experiment, the code was deliberately reconfigured with an incorrect address (\texttt{0x28} instead of \texttt{0x27}).

\begin{table}[h!]
\centering
\small
\caption{System Behavior Under Address Mismatch (\texttt{0x28})}
\vspace{4pt}
\begin{tabularx}{\textwidth}{l p{4.5cm} p{4.5cm} X}
\toprule
\textbf{Platform} & \textbf{Function Call} & \textbf{Returned Status Code} & \textbf{Physical Observation \& Bus State} \\
\midrule
\textbf{STM32 HAL} 
    & \texttt{HAL\_I2C\_Master\_Transmit} 
    & \texttt{HAL\_ERROR} (\texttt{0x01}) \newline Error Flag: \texttt{HAL\_I2C\_ERROR\_AF} (\texttt{0x04}) 
    & The LCD remains completely blank or shows dark contrast blocks. No ACK is generated on the 9th clock cycle of the address byte. The transmission aborts immediately. \\
\midrule
\textbf{Arduino Wire} 
    & \texttt{Wire.endTransmission()} 
    & Return code \textbf{2} (\texttt{NACK on address}) 
    & Text fails to appear. The driver receives a NACK after sending the invalid address \texttt{0x28} and refuses to stream subsequent character bytes. \\
\bottomrule
\end{tabularx}
\end{table}

% =========================================================================
\section{Timing \& Wire Throughput Analysis (Task Question 4)}
% =========================================================================

\begin{questionbox}{4}
\textbf{The STM32 driver sends four bytes per LCD character in one transaction. How many bits are on the wire per character at $100\text{ kHz}$ and how long does it take? How many characters per second can the display be refreshed at most?}
\end{questionbox}

\subsubsection*{1. Bit Count on the Physical Wire per Character}
In our driver, one LCD character is transmitted as a single atomic transaction consisting of:
\begin{itemize}
    \item 1 START condition ($S$)
    \item 1 Slave Address Byte (\texttt{0x4E}): 8 data bits + 1 ACK bit = \textbf{9 clock cycles}
    \item 4 Payload Bytes (Upper nibble with $E=1$, Upper nibble with $E=0$, Lower nibble with $E=1$, Lower nibble with $E=0$):
    \begin{equation}
    4\text{ data bytes} \times (8\text{ data bits} + 1\text{ ACK bit}) = 4 \times 9 = \mathbf{36\text{ clock cycles}}
    \end{equation}
    \item 1 STOP condition ($P$)
\end{itemize}
Total clock periods on the wire per character:
\begin{equation}
N_{\text{clocks}} = 9\text{ (address)} + 36\text{ (data)} = \mathbf{45\text{ SCL clock periods}}
\end{equation}

\subsubsection*{2. Transmission Duration at $100\text{ kHz}$ ($T_{\text{SCL}} = 10.0\,\mu\text{s}$)}
\begin{equation}
t_{\text{char}} = 45\text{ clocks} \times 10.0\,\mu\text{s/clock} = \mathbf{450\,\mu\text{s}} \quad (0.45\text{ ms})
\end{equation}
*(If calculating pure data payload without address re-transmission in a streaming transaction: $36 \times 10\,\mu\text{s} = 360\,\mu\text{s}$)*.

\subsubsection*{3. Maximum Characters Refreshed Per Second}
The theoretical maximum continuous character transfer rate across the bus:
\begin{equation}
\text{Throughput} = \frac{1}{t_{\text{char}}} = \frac{1}{450\times 10^{-6}\text{ s}} \approx \mathbf{2222.2\text{ characters/second}}
\end{equation}
\textbf{Full Screen Refresh Rate:}  
A $16\times 2$ LCD has 32 visible character cells. Refreshing all 32 cells requires:
\begin{equation}
t_{\text{screen}} = 32 \times 450\,\mu\text{s} = 14.4\text{ ms} \implies \text{Max Frame Rate} = \frac{1}{0.0144\text{ s}} \approx \mathbf{69.4\text{ frames/second}}
\end{equation}
This demonstrates that even at standard $100\text{ kHz}$ speeds, the $\text{I}^2\text{C}$ bus can refresh the entire display at $\approx 70\text{ Hz}$, far surpassing human visual persistence.

\end{document}
